\chapter{Introduction}
\label{chap:introduction}

\section{Background and Rationale}
\label{sec:background}

Video has become the dominant form of content on the Internet, and platforms such as
YouTube, Vimeo, and TikTok have made uploading and sharing video an everyday activity. Behind
that apparent simplicity lies a demanding engineering problem. When a user uploads a source
file, that file may use any container format, any codec, and any resolution, yet the platform
must deliver a stream that plays smoothly on every device and adapts to whatever bandwidth the
viewer happens to have at that moment.

The industry solution to this problem is HTTP Live Streaming (HLS), a protocol that splits a
video into short segments, typically between two and ten seconds long, and produces several
quality renditions of the same content. A playlist file written in the \texttt{.m3u8} format
enumerates the available renditions, and the player selects among them in real time according
to measured network throughput. This mechanism, known as Adaptive Bitrate Streaming (ABR),
is what allows playback to degrade gracefully rather than stall when a connection weakens.

Producing those renditions is computationally expensive, and this is where the architectural
difficulty arises. Transcoding is a bursty workload: it consumes a large amount of CPU while a
video is being processed and nothing at all when no one is uploading. A platform that keeps a
dedicated server cluster or a Kubernetes system running continuously simply to wait for work
therefore pays for a substantial amount of idle capacity. The alternative failure mode is
equally undesirable: a system provisioned only for the average case cannot absorb a sudden
burst when many users upload at the same time, and uploads either queue for a long time or are
dropped altogether.

This tension between idle cost and burst capacity motivates the present project. Rather than
maintaining permanently allocated compute, the system proposed here starts a container only
when there is a video to process and destroys it immediately afterwards, so that cost accrues
only for the seconds during which transcoding actually runs.

\section{Project Objectives}
\label{sec:objectives}

The project aims to build a complete video sharing platform for end users, together with the
automated transcoding subsystem that supports it. Specifically, the objectives are:

\begin{enumerate}
    \item To build a web platform that allows users to register an account, upload videos,
          watch them at multiple quality levels, and share them through either public or
          private links.
    \item To automate the entire processing pipeline, from the moment a source file lands in
          Amazon S3 through transcoding, storage, and delivery via a content delivery network,
          until the video becomes ready for playback in a browser.
    \item To perform transcoding with FFmpeg inside a Docker container running on AWS Batch and
          AWS Fargate, such that containers are created on demand and terminate automatically,
          thereby eliminating the cost of idle resources.
    \item To automate software deployment and infrastructure management through a CI/CD pipeline
          and Infrastructure as Code.
    \item To evaluate the resulting system empirically in terms of quality of experience,
          transcoding performance, scalability under load, and operating cost.
\end{enumerate}

\section{Scope and Deliverables}
\label{sec:scope}

The scope of this project covers the full vertical slice of a video sharing service: the
browser-based user interface, the backend API, the transcoding engine, the cloud
infrastructure that hosts them, and the automation that deploys them.

Several areas are deliberately placed outside the scope. The project does not implement
digital rights management systems such as Widevine or FairPlay, which would be required to
prevent an authorised viewer from redistributing content they can legitimately access. It does
not implement live streaming, since the architecture is designed around video on demand. It
also does not attempt per-title bitrate ladder optimisation, using instead a fixed ladder of
three renditions; the rationale and the trade-off involved are discussed in
Chapter~\ref{chap:conclusion}.

\section{Research Methodology}
\label{sec:methodology}

The work combines a review of existing practice with an iterative build-and-measure cycle.

The theoretical component established how the problem is solved at industrial scale before any
code was written. Published engineering accounts from Netflix describe how a large platform
decomposes video processing into microservices and how a bitrate ladder is optimised per title,
and the AWS Video on Demand reference architecture describes the managed-service equivalent. The
protocol itself was taken from its specification, RFC 8216, and the encoding parameters from the
FFmpeg HLS muxer documentation, rather than from secondary sources. Peer-reviewed studies of
serverless video processing supplied both the motivation for the architecture and the measurement
vocabulary used later in the evaluation, and the multi-threaded H.264 literature supplied the
expected bound against which the vCPU scaling results of Chapter~\ref{chap:evaluation} are
interpreted.

The practical component followed an incremental development model. Each subsystem, comprising the
backend API, the transcoding container, the frontend, and the cloud infrastructure, was built and
integrated in turn, with the whole environment described as Terraform code so that every change to
it is reviewable and reproducible rather than applied by hand. Correctness was checked by an
automated test suite executed on every push through a continuous integration pipeline, which also
runs dependency, container and infrastructure security scans.

Evaluation was empirical rather than analytical. Every performance figure reported in
Chapter~\ref{chap:evaluation} was produced by a measurement script committed alongside the
source code, executed against the deployed system rather than a local approximation of it, with
the raw output stored as JSON under \path{docs/results/} so that any table in this report can be
traced back to the run that produced it. Where a result could plausibly be an artefact of the
measurement tool, it was reproduced with a second, independently implemented tool before being
reported.

\section{Practical Significance}
\label{sec:significance}

The immediate beneficiaries of this design are small teams and individual creators who need a
video platform but cannot justify a permanently running transcoding cluster. The cost comparison
in Section~\ref{sec:cost-analysis} estimates that processing three hundred videos a month costs
roughly \$1.51 on the architecture built here against \$133.35 on an always-on instance sized for
the same work, because compute exists only while a job runs. That difference is what makes video
hosting feasible for a service whose upload volume is intermittent, which describes most
platforms before they become large.

The design also addresses a requirement that is easy to state and easy to implement incorrectly.
Marking a video private is meaningless if the underlying storage remains publicly addressable, and
Chapter~\ref{chap:design} shows that restricting the API alone does not achieve it. The two-layer
scheme used here, keeping the bucket private and enforcing authorisation at the content delivery
network, is directly reusable by any application that serves user-owned media from object storage.

Beyond the product itself, the project serves as a worked reference for an event-driven serverless
pipeline on AWS, covering the parts that are usually omitted from tutorials: how duplicate message
delivery is made harmless, how the infrastructure is expressed as reviewable code across two
environments, how security scanning is placed in the path of a merge rather than beside it, and
how the resulting system is measured once it is running. Each of these is documented here in
enough detail to be repeated.

\section{Main Contributions}
\label{sec:contributions}

The principal contribution of this work is an end-to-end demonstration that a serverless
container architecture can serve as the transcoding backbone of a video platform, together
with an empirical assessment of what that choice costs and what it saves.

More concretely, the work contributes an event-driven pipeline in which an upload event
propagates through Amazon S3, Amazon SQS, and AWS Lambda to AWS Batch without any polling
service in between; a containerised FFmpeg worker whose conditional database writes make a
duplicate delivery harmless, so that exactly-once \emph{effects} are obtained without relying on
exactly-once delivery; a two-layer content protection scheme combining Origin Access Control with
CloudFront Signed Cookies so that private videos remain inaccessible even to a viewer holding
the exact manifest URL; a complete Infrastructure as Code description of the AWS environment
in eleven Terraform modules; and a DevSecOps pipeline whose security gate genuinely blocks
merges rather than merely reporting findings.

\section{Structure of the Report}
\label{sec:structure}

The remainder of this report is organised as follows. Chapter~\ref{chap:literature} reviews
how comparable platforms solve the transcoding problem and justifies the technology choices
made here. Chapter~\ref{chap:requirements} analyses the functional and non-functional
requirements of the system and presents the use case model. Chapter~\ref{chap:design}
describes the architecture, the data model, the content protection scheme, and the interface
design. Chapter~\ref{chap:implementation} details the implementation of each subsystem, the
infrastructure automation and continuous integration, the difficulties encountered during
development and how they were resolved, and the automated test suite.
Chapter~\ref{chap:evaluation} presents the deployed system and its user interface, then reports
the experimental results for delivery latency, transcoding performance, behaviour under load, and
operating cost. Chapter~\ref{chap:conclusion} summarises what was achieved, states the strengths
and limitations of the delivered system, and outlines directions for future work.
